\chapter{Related Work}
\label{ch:related}

This chapter reviews the prior work that \textsc{EdgeGuard} builds on or competes with. It is
organised by threat: fault attacks on weights (\cref{sec:rw:bfa}) and defences against them
(\cref{sec:rw:bfadef}); adversarial inputs and feature-space detection (\cref{sec:rw:adv});
extraction of models and training data (\cref{sec:rw:extract}); supply-chain manipulation of
weights (\cref{sec:rw:supply}); and TEE-based protection (\cref{sec:rw:tee}).
\Cref{sec:rw:gap} positions this project against them.

% ---------------------------------------------------------------------------
\section{Fault attacks on model weights}
\label{sec:rw:bfa}

\paragraph{Rowhammer.} Kim et al.\ showed that repeatedly activating a DRAM row can flip bits
in physically adjacent rows without accessing them \cite{kim2014flipping}. The effect gives
unprivileged software a write primitive into memory it does not own, and it has since been
exploited from cloud tenants, browsers and mobile applications.

\paragraph{Sensitivity of networks to single flips.} Neural networks tolerate random noise in
their parameters well, which long suggested that hardware faults were a minor concern. Hong et
al.\ challenged this view \cite{hong2019terminal}. Across 19 models covering six architectures
and three image tasks, most models had at least one parameter whose single bit flip caused an
accuracy loss of more than 90\%, and an estimated 40--50\% of parameters could cause a drop of
more than 10\% when flipped individually. Their analysis covered floating-point models, where
flipping a high exponent bit is catastrophic. The same work demonstrated a practical Rowhammer
attack in a simulated machine-learning-as-a-service setting.

\paragraph{Attacks on quantized models.} Quantization was initially thought to remove the
exponent-bit vulnerability. Rakin, He and Fan's bit-flip attack (BFA) showed otherwise
\cite{rakin2019bfa}. It uses gradient information to search progressively for the bits whose
flipping most increases the loss, and it degrades 8-bit quantized ResNet models to near-random
accuracy with a few tens of flips out of tens of millions of weight bits. DeepHammer by Yao,
Rakin and Fan turned this search into an end-to-end hardware attack \cite{yao2020deephammer}.
It constrains the search to bits that are actually flippable on the victim's DRAM, uses memory
massaging to place victim weight pages on vulnerable rows, and degraded the accuracy of 12
architectures across four datasets to the level of random guessing within minutes. DeepHammer
also reports that randomly flipping even 100 bits in an 8-bit ResNet-20 barely changes accuracy,
so effective attacks are necessarily targeted. This observation shapes the attack simulators in
\cref{sec:impl:attacks}.

% ---------------------------------------------------------------------------
\section{Defences against weight fault attacks}
\label{sec:rw:bfadef}

Defences fall into three groups.

\paragraph{Making the model harder to attack.} He et al.\ studied binarization-aware training
and piece-wise weight clustering as ways to raise the number of flips an attacker needs
\cite{he2020defending}. Li et al.\ proposed reconstructing weights at inference time so that the
effect of a flipped bit is spread over a group, which maintained 60\% ImageNet accuracy after
five iterations of a greedy progressive attack, compared with below 1\% for the undefended
baseline \cite{li2020defending}. These approaches change the model or its training and do not
report whether an attack occurred.

\paragraph{Detecting modification with signatures.} RADAR by Li et al.\ groups interleaved
weights, derives a 2-bit addition-checksum signature per group, and compares it at run time with
a securely stored reference \cite{li2021radar}. It detects about 9.6 out of 10 injected flips on
ResNet-18 and zeroes flagged groups to recover accuracy, adding under 1\% to inference time in a
gem5 simulation. HASHTAG by Javaheripi and Koushanfar identifies the layers most sensitive to
fault injection and hashes their weights with a low-collision (Pearson) hash, comparing against
ground-truth signatures during inference \cite{javaheripi2021hashtag}. ModelShield by Guo et al.\
performs real-time hash verification of weights as an extension to existing frameworks
\cite{guo2021modelshield}. These works establish that hashing weights at run time is effective
against BFA, and \textsc{EdgeGuard}'s attestation component belongs to the same family. It
differs in using a collision-resistant hash over every tensor, a Merkle tree for localisation,
and a signature on the root. Lightweight checksums can be defeated by an attacker who knows the
checksum and flips compensating bits. A cryptographic hash cannot be defeated this way, at the
price of higher cost, which \cref{sec:eval:merkle} measures.

\paragraph{Luring the attacker.} NeuroPots by Liu et al.\ embeds ``honey neurons'' into the
model as deliberately vulnerable points \cite{liu2023neuropots}. Attack searches tend to
concentrate their flips there, so a checksum over the small trapdoor region is enough to detect
the attack and recover the model. The approach reduces the region that must be checked, which
directly addresses the cost problem.

% ---------------------------------------------------------------------------
\section{Adversarial inputs and feature-space detection}
\label{sec:rw:adv}

Goodfellow et al.\ explained adversarial examples by the locally linear behaviour of networks
and introduced the fast gradient sign method \cite{goodfellow2015explaining}. Madry et al.\
formulated robustness as a saddle-point problem and used projected gradient descent both to
attack and to train robust models \cite{madry2018towards}.

Detection, as opposed to robust training, often works in feature space. Lee et al.\ fit a
class-conditional Gaussian with a shared covariance to the features of each layer of a trained
classifier \cite{lee2018simple}. They score a test input by its Mahalanobis distance to the
closest class mean and combine scores across layers. The method detects both out-of-distribution
samples and adversarial examples without retraining the classifier. \textsc{EdgeGuard}'s
monitor uses the same statistic in a simpler, class-agnostic form: one Gaussian per monitored
layer, fitted on INT8 activations from clean inputs. Being class-agnostic makes it applicable to
generative models, where there are no class labels. It is also less discriminative, because the
clean distribution it models is broader.

% ---------------------------------------------------------------------------
\section{Extraction of models and training data}
\label{sec:rw:extract}

\paragraph{Model extraction.} Tram\`er et al.\ showed that models exposed through prediction
APIs can be replicated with a number of queries close to the number of parameters for simple
model classes \cite{tramer2016stealing}. With physical access, extraction is more direct.
Halderman et al.\ showed that DRAM retains its contents for seconds to minutes after power
loss, longer when cooled, which allows memory contents to be recovered by rebooting into a
custom image \cite{halderman2008lest}. The cold-boot channel model in
\cref{sec:impl:attacks} simulates such a recovery with a small rate of single-bit decay errors.

\paragraph{Training-data extraction.} Carlini et al.\ generated large numbers of samples from
GPT-2 and ranked them with membership-inference metrics. Hundreds of the top-ranked samples
were verbatim training sequences, including names, phone numbers and addresses
\cite{carlini2021extracting}. Nasr et al.\ extended extraction to open and closed production
models and showed that aligned chat models can be prompted to emit training data at much
higher rates than in normal use \cite{nasr2023scalable}. Dai, Lu and Zhou introduced the
activation inversion attack \cite{dai2025activation}. In decentralised (pipeline-parallel)
training, an attacker who observes the activations passed between stages trains a model on a
shadow dataset of text--activation pairs and uses it to reconstruct the training text. That
work concerns activations transmitted during training rather than inference on a device, but it
establishes that intermediate activations carry recoverable information about the input text.

The extraction-probe simulator in \textsc{EdgeGuard} does not implement any of these attacks.
It models a hypothesis: queries that make a model regurgitate memorised content produce
abnormally confident, low-variance late-layer activations. This hypothesis is plausible, but it
has not been validated on a real model in this project, and \cref{ch:discussion} treats it
accordingly.

% ---------------------------------------------------------------------------
\section{Supply-chain manipulation of weights}
\label{sec:rw:supply}

Practitioners routinely download pretrained weights from public repositories and fine-tune
them. Feng and Tram\`er showed that an attacker who tampers with the pretrained weights can
plant ``data traps'' that record individual fine-tuning examples into the weights, allowing the
attacker to reconstruct them later \cite{feng2024privacy}. The traps can also yield perfect
membership inference and defeat loose differential-privacy guarantees. Wen et al.\
independently showed that poisoned pretrained models can amplify membership-inference leakage
from fine-tuning data \cite{wen2024privacy}. Both results make the provenance of weights a
security property. A signed attestation of the weights shipped to a device cannot by itself
reveal a backdoor planted before signing. It can, however, ensure that the weights on the device
are the ones the vendor signed, so that tampering after signing is detected.

% ---------------------------------------------------------------------------
\section{Trusted execution for model protection}
\label{sec:rw:tee}

TEEs protect code and data from the rest of the system but have limited memory and no
accelerator access. Amulet by Mao et al.\ obfuscates the weights inside a TEE, stores the
transformed model in untrusted memory, and runs it on an external accelerator with only two TEE
interactions per inference \cite{mao2025amulet}. It reports near-identical accuracy with latency
within $4.8\times$ of unprotected GPU inference across models from ResNet-18 to GPT-2.
Amulet targets \emph{confidentiality} of the weights against extraction, whereas
\textsc{EdgeGuard} targets \emph{integrity} and \emph{runtime anomaly detection}. The two are
complementary. In particular, a TEE is the natural home for \textsc{EdgeGuard}'s verification
key and attestation logic.

% ---------------------------------------------------------------------------
\section{Positioning of this work}
\label{sec:rw:gap}

\Cref{tab:related} compares the systems discussed above along the dimensions relevant to this
project. Each entry reflects what the cited paper reports. A circle (\pmark) marks partial
coverage.

\begin{table}[t]
\centering
\caption{Coverage of related defences. \cmark: addressed; \pmark: partially; \xmark: not
addressed in the cited work. ``Localises'' means the defence identifies where the modification
occurred.}
\label{tab:related}
\small
\resizebox{\textwidth}{!}{\begin{tabular}{lcccccc}
\toprule
\textbf{System} & \textbf{Detects} & \textbf{Crypto.} & \textbf{Localises} & \textbf{Input} &
\textbf{Labels} & \textbf{Evaluated} \\
 & \textbf{weight mod.} & \textbf{strength} & & \textbf{anomalies} & \textbf{threat} & \textbf{on LLMs} \\
\midrule
Weight reconstr.\ \cite{li2020defending} & \xmark & \xmark & \xmark & \xmark & \xmark & \xmark \\
RADAR \cite{li2021radar}                 & \cmark & \xmark & \cmark (group) & \xmark & \xmark & \xmark \\
HASHTAG \cite{javaheripi2021hashtag}     & \cmark & \xmark & \pmark (layer) & \xmark & \xmark & \xmark \\
ModelShield \cite{guo2021modelshield}    & \cmark & \pmark & \pmark & \xmark & \xmark & \xmark \\
NeuroPots \cite{liu2023neuropots}        & \cmark & \xmark & \cmark (trapdoor) & \xmark & \xmark & \xmark \\
Mahalanobis det.\ \cite{lee2018simple}   & \xmark & \xmark & \xmark & \cmark & \xmark & \xmark \\
Amulet \cite{mao2025amulet}              & \xmark & \xmark & \xmark & \xmark & \xmark & \pmark (GPT-2) \\
\midrule
\textsc{EdgeGuard} (this work)           & \cmark & \cmark & \cmark (tensor) & \cmark & \cmark & \pmark (stub) \\
\bottomrule
\end{tabular}}
\end{table}

Two points in \cref{tab:related} need qualification. First, \textsc{EdgeGuard}'s ``evaluated on
LLMs'' is marked partial because only structural stubs are used. Second, the table records
coverage, not effectiveness. RADAR, HASHTAG and NeuroPots were evaluated on real models under
real attacks, whereas \textsc{EdgeGuard}'s activation results come from simulation. The defensible
claim of this project is narrower than ``a better defence''. It is a working integration of
cryptographic attestation with INT8 activation monitoring and threat labelling, with a measured
characterisation of what each part contributes and what it costs.
